\title{Controlling Text-to-Image Diffusion by Orthogonal Finetuning}

\begin{abstract}
Large-scale text-to-image diffusion models have demonstrated remarkable proficiency in synthesizing photorealistic visuals from textual descriptions. A central, unresolved question is how to effectively direct or modulate these models for various downstream applications. In response, we propose Orthogonal Finetuning~(OFT), a theoretically grounded method for adapting text-to-image diffusion models to specialized tasks. Distinct from prior approaches, OFT can provably retain hyperspherical energy, which encapsulates the relationships between neurons on the unit hypersphere. Preserving this attribute proves vital for maintaining the semantic image generation abilities of text-to-image diffusion models. For enhanced stability during finetuning, we further introduce Constrained Orthogonal Finetuning (COFT), incorporating an extra constraint on the hypersphere radius. Our study focuses on two representative finetuning scenarios: subject-driven generation, which involves producing subject-specific images given limited subject examples and a guiding text prompt, and controllable generation, where models are equipped to integrate supplementary control signals. Experimental results reveal that our OFT framework surpasses existing approaches regarding both image quality and convergence speed.
\end{abstract}

\section{Introduction}

State-of-the-art text-to-image diffusion models~\cite{saharia2022photorealistic,ramesh2022hierarchical,rombach2022high} excel in generating high-resolution images based on textual input. Nevertheless, the guidance derived from text alone is often imprecise and insufficient to achieve detailed or highly controlled image outcomes. To mitigate these limitations, this paper investigates two principal text-to-image generation challenges:

\begin{itemize}[leftmargin=*,nosep]

    \item \textbf{Subject-driven generation}~\cite{ruiz2023dreambooth}: With access to a small set of images depicting a subject, this task seeks to generate new images of the same subject in novel contexts, leveraging textual instructions.
    \item \textbf{Controllable generation}~\cite{zhang2023adding,mou2023t2i}: Here, the model is provided with additional control signals (\emph{e.g.}, canny edge maps, segmentation masks) alongside text prompts, aiming to produce images that adhere to these control inputs.
\end{itemize}

Both of these tasks ultimately revolve around the challenge of finetuning text-to-image diffusion models in a way that maintains the generative quality acquired during pretraining. The main requirements for a robust finetuning technique can be distilled as: (1) \emph{training efficiency}, which refers to minimizing both the number of trainable parameters and the training duration, and (2) \emph{generalizability preservation}, which entails retaining the pretrained model’s capacity for generating high-quality and diverse outputs. Existing approaches to finetuning typically either involve directly modifying the model's weights using a low learning rate (\emph{e.g.}, \cite{ruiz2023dreambooth}), or introducing additional components with reparameterized weights (\emph{e.g.}, \cite{hulora2022,zhang2023adding}). Despite their straightforwardness, these strategies fall short of ensuring that the generative abilities conferred by pretraining are preserved. Furthermore, there remains a lack of theoretical grounding for developing effective finetuning methods, as well as for determining crucial hyperparameters like the number of training epochs. One core issue lies in the absence of a concrete metric for assessing how well the pretrained generative capabilities are retained. Many existing approaches operate under the implicit assumption that maintaining a small Euclidean distance between the finetuned and pretrained models corresponds to better retention of original generative abilities, and consequently employ very conservative learning rates. Although this assumption might occasionally prove valid, we contend that Euclidean proximity by itself is insufficient for encapsulating semantic consistency. Therefore, adopting a more structurally informed metric to gauge the distinction between the finetuned and pretrained models holds promise for enhancing both the maintenance of pretrained performance and the reliability of the finetuning process.

Building on prior empirical findings that hyperspherical similarity effectively captures semantic relationships~\cite{liu2017deep,liu2018decoupled,chen2020angular}, our approach leverages hyperspherical energy~\cite{liu2018learning} as a means to represent the pairwise interactions among neurons. Specifically, hyperspherical energy is computed by summing the hyperspherical similarity (\emph{e.g.}, cosine similarity) for all pairs of neurons within a layer, thus reflecting the degree of uniformity in their arrangement on the unit hypersphere~\cite{liu2021learning}. We propose that optimal finetuning should strive to minimize the change in hyperspherical energy relative to the pretrained network. Although a straightforward method might introduce a regularization term to preserve hyperspherical energy throughout finetuning, such an approach does not guarantee minimal energy variation. Instead, we utilize an invariance property of hyperspherical energy, which states that pairwise hyperspherical similarity is maintained under a common orthogonal transformation applied to all neurons. Drawing on this invariance, we introduce Orthogonal Finetuning~(OFT), which adapts large-scale text-to-image diffusion models for new tasks while preserving their hyperspherical energy. The key idea is to learn an orthogonal transformation shared across layers that preserves pairwise angular relationships among neurons. This process can be interpreted as redefining the canonical coordinate system for neurons within a layer. Recognizing that reduced Euclidean distance between the finetuned and pretrained models correlates with improved retention of pretrained performance, we further introduce a refined method, Constrained Orthogonal Finetuning~(COFT), which restricts the finetuned parameters to lie on a hypersphere of fixed radius centered at the original pretrained neurons.

The rationale behind employing orthogonal transformation in neuron finetuning partially draws on the analogy with the 2D Fourier transform, wherein an image is represented through its magnitude and phase spectra. The phase spectrum encodes the angular relationship between the input and the basis functions, and is chiefly responsible for retaining the semantic information. Notably, even if the magnitude spectrum is randomized, reconstructing an image using the original phase spectrum is sufficient to largely preserve the image’s semantics. This observation indicates that adjusting neuron directions fundamentally underpins semantic alteration in generated images—an objective shared by both subject-driven and controllable generative tasks. Nonetheless, unconstrained modification of neuron directions risks compromising the generative abilities instilled during pretraining due to the resulting large degrees of freedom. To address this, we advocate for maintaining the pairwise angles between neurons, based on the premise that these angular relationships encapsulate critical knowledge within neural networks. Building on this insight, we propose learning a layer-wise, shared orthogonal transformation for neurons, thereby ensuring invariance of the hyperspherical energy.

Our methodology is further motivated by prior work in orthogonal over-parameterized training~\cite{liu2021orthogonal}, which constructs classification networks from the ground up through orthogonal transformation of randomly initialized parameters. In such cases, the expected hyperspherical energy of random initialization is provably low, and the objective of~\cite{liu2021orthogonal} is to maintain this low hyperspherical energy throughout training—a property correlated with improved classification generalization~\cite{liu2018learning,lin2020regularizing}. The findings in~\cite{liu2021orthogonal} demonstrate that orthogonal transformations possess sufficient flexibility for training effective classification models. In contrast, our focus is on adapting text-to-image diffusion models for enhanced controllability and superior generative capability in downstream applications. It is important to distinguish OFT from~\cite{liu2021orthogonal} in two principal respects. First, whereas~\cite{liu2021orthogonal} is geared towards minimizing hyperspherical energy, OFT strives to retain the same hyperspherical energy as the pretrained model, ensuring the structural integrity of the pretrained representation remains unaffected during finetuning. Minimizing hyperspherical energy in diffusion model finetuning can be detrimental, as it might disrupt essential semantic structures. Second, OFT introduces mechanisms to limit divergence from the pretrained parameters, resulting in a constrained formulation, unlike~\cite{liu2021orthogonal}, which omits such restrictions. Ultimately, finetuning entails balancing adaptability with preservation of model stability, and we posit that our OFT approach provides a robust solution to this trade-off. Our main contributions are outlined below:

\begin{itemize}[leftmargin=*,nosep]

    \item We introduce a new finetuning strategy—Orthogonal Finetuning—that enhances the controllability of text-to-image diffusion models. To increase robustness, we further present a constrained variant, restricting the angular shift relative to the original pretrained weights.
    \item Unlike previous finetuning approaches, OFT adjusts model parameters in a manner that rigorously maintains the hyperspherical energy, a property we empirically associate with the retention of the semantic generative abilities inherent to the pretrained model.
    \item We implement OFT in the contexts of subject-driven generation and controllable generation. Through extensive experimental evaluation, we demonstrate marked improvements over existing approaches regarding image generation fidelity, training convergence rates, and overall stability during finetuning. Additionally, OFT is more sample-efficient, achieving reliable convergence with significantly fewer training samples and epochs.
    \item For the controllable generation scenario, we propose a novel control consistency metric to quantify control efficacy. This approach involves extracting the control signal from each synthesized image and comparing it to the original input control. Our metric provides additional evidence for the strong controllability enabled by OFT.
\end{itemize}

\section{Related Work}

\textbf{Text-to-image diffusion models}. Recent years have witnessed significant advances in text-to-image synthesis~\cite{saharia2022photorealistic,ramesh2022hierarchical,rombach2022high,gu2022vector,nichol2022glide}, primarily fueled by the progress in diffusion-based generative models~\cite{ho2020denoising,song2021score,song2021denoising,dhariwal2021diffusion} as well as advancements in vision-language representation learning~\cite{su2020vlbert,lu2019vilbert,radford2021learning,wang2021simvlm,li2022blip,li2023blip,alayrac2022flamingo,schuhmann2022laion}. Models such as GLIDE~\cite{nichol2022glide} and Imagen~\cite{saharia2022photorealistic} operate directly in pixel space: GLIDE learns both the text encoder and diffusion prior via paired data, whereas Imagen employs a fixed, pretrained text encoder throughout. Alternatively, Stable Diffusion~\cite{rombach2022high} and DALL-E2~\cite{ramesh2022hierarchical} are trained in latent space: Stable Diffusion leverages VQ-GAN~\cite{esser2021taming} for constructing a visual codebook, while DALL-E2 utilizes CLIP~\cite{radford2021learning} to create a unified latent embedding for text and images. In addition to diffusion methods, text-to-image generation has been pursued using generative adversarial networks~\cite{reed2016generative,zhang2017stackgan,xu2018attngan,li2019controllable} and autoregressive techniques~\cite{ramesh2021zero,ding2021cogview,wu2022nuwa,yuscaling}. OFT is designed as a model-agnostic finetuning mechanism, making it compatible with any text-to-image diffusion framework.

\textbf{Subject-driven generation}. Addressing undesired subject alterations, \cite{avrahami2022blended,nichol2022glide} incorporate user-provided masks as additional conditioning. Inversion-based strategies~\cite{choi2021ilvr,dhariwal2021diffusion,ramesh2022hierarchical,gal2022image} permit contextual changes while preserving the subject, and \cite{hertz2022prompt} can edit content at both local and global scales without needing an explicit mask. However, these approaches frequently struggle to faithfully maintain identity-specific features. The Pivotal Tuning method~\cite{roich2022pivotal} adapts a generator around an initial latent inversion, supplemented by regularization for identity conservation. Similarly, \cite{nitzan2022mystyle} develops a personalized facial generative prior from a set of individual images, while \cite{casanova2021instance} enables the synthesis of variations on a given instance—though potentially at the cost of some instance-level specificity. DreamBooth~\cite{ruiz2023dreambooth} finetunes text-to-image diffusion models by leveraging a custom token and several subject images, utilizing reconstruction and class-specific prior losses for supervision. OFT builds on the DreamBooth protocol, but diverges by employing orthogonal updates rather than standard gradient-based finetuning with small learning rates.

\textbf{Controllable generation}. Image-to-image translation can be framed as a controllable generation problem. Traditionally, this challenge has been tackled with conditional generative adversarial networks~\cite{isola2017image,zhu2017toward,wang2018high,park2019semantic,choi2018stargan}. Recent advancements also see diffusion models being leveraged for the task~\cite{saharia2022palette,voynov2022sketch,wang2022pretraining}. ControlNet~\cite{zhang2023adding} represents a significant step by adapting a pretrained diffusion model through finetuning to accommodate new control signals, achieving state-of-the-art performance in controllable generation. Similarly, T2I-Adapter~\cite{mou2023t2i} refines a pretrained diffusion model to improve controllability over generated outputs. Building upon the experimental protocols established in \cite{zhang2023adding,mou2023t2i}, we utilize OFT to finetune diffusion models. OFT consistently results in superior controllability, requiring less training data and fewer additional finetuning parameters. Importantly, OFT does not incur extra computational cost at inference time.

\textbf{Model finetuning}. Tailoring large pretrained models for downstream tasks has become increasingly prevalent~\cite{devlin2018bert,brown2020language,he2022masked}. Within this paradigm, adaptation techniques (\emph{e.g.}, \cite{hulora2022,houlsby2019parameter,pfeiffer2020adapterfusion}) have been extensively explored, especially in natural language processing. Of particular relevance to OFT is LoRA~\cite{hulora2022}, which employs a low-rank approximation for weight updates in the finetuning phase. However, unlike LoRA, OFT utilizes a shared orthogonal transformation per layer to modify neuron weights multiplicatively, while theoretically preserving the pairwise angles among neurons within the layer, leading to increased stability.

\section{Orthogonal Finetuning}

\subsection{Motivation for Orthogonal Transformation}

We begin by elaborating on the rationale behind employing orthogonal transformations when finetuning text-to-image diffusion models. Specifically, the discussion is separated into two parts: (1) the motivation for finetuning the direction (or angle) of neurons, and (2) the justification for using orthogonal transformations for this purpose.

To address the first question, we build on prior empirical findings in \cite{liu2018decoupled,chen2020angular} that suggest angular differences between features are effective at representing the semantic disparity. SphereNet~\cite{liu2017deep} provides evidence that enforcing normalization of neural activations onto the unit hypersphere can yield models with competitive expressiveness and superior generalization, which indicates that the directional components of neurons capture essential data attributes. To highlight the significance of neuron orientation, we present a controlled experiment depicted in Figure~\ref{fig:toy_ae}. Here, a typical convolutional autoencoder is trained on images of flowers. During training, the usual inner product forms the feature map ($z$ is the output for a kernel $\bm{w}$ on an input $\bm{x}$ within a sliding window). When evaluating, we apply three methods for producing the feature map: (a) the training-phase inner product, (b) solely considering magnitude, and (c) isolating angular information. The outcomes in Figure~\ref{fig:toy_ae} reveal that the angular component of the neural activations nearly reconstructs the original images perfectly, whereas magnitude information does not retain useful content. Notably, the network was trained using only the inner product, without incorporating the cosine-based (angular) activation (c) during training. This indicates that semantic details are predominantly encoded in neuron directions rather than their magnitudes, suggesting that modifying neuron orientation is likely the most effective strategy for adjusting semantic image content.

With respect to the second question, a direct approach to adjusting neuron directions is to apply collective rotations or reflections to all neurons within a layer, which corresponds to an orthogonal transformation. While more intricate angular transformations—such as rotating individual neurons by distinct angles—are possible, our experiments indicate that orthogonal transformations strike an ideal compromise between adaptability and structural regularity. In addition, the findings in \cite{liu2021orthogonal} demonstrate the expressiveness of orthogonal transformations in neural network training. To reinforce our point, we conduct an experiment examining the regularizing effect of enforcing orthogonality. Utilizing the experimental framework of ControlNet~\cite{zhang2023adding}, we report results in Figure~\ref{fig:ortho}. Our standard OFT approach demonstrates consistent behavior and precise controllability after 20 epochs of training, while omitting the orthogonality condition results in the failure to generate realistic or controlled images. This provides strong empirical support for the critical role of orthogonality constraints in OFT.

\subsection{General Framework}

Standard finetuning approaches generally employ gradient descent with a reduced learning rate, updating either the entire model or selected layers. Using a low learning rate acts as an implicit mechanism to restrict the finetuned model's departure from its pretrained configuration; fundamentally, standard finetuning encourages minimizing $\thickmuskip=2mu \medmuskip=2mu \| \bm{M} - \bm{M}^0\|$, where $\bm{M}$ and $\bm{M}^0$ represent the parameters of the finetuned and pretrained models, respectively. While this implicit regularization is intuitively reasonable, it may still lack sufficient rigidity, particularly for large-scale models. LoRA tackles this by enforcing an explicit low-rank constraint on the model update: $\thickmuskip=2mu \medmuskip=2mu \text{rank}(\bm{M} - \bm{M}^0)=r'$, with $r'$ set to a small value. In contrast, OFT imposes a restriction on neuron-wise pairwise similarity, formalized as $\thickmuskip=2mu \medmuskip=2mu \| \text{HE}(\bm{M})-\text{HE}(\bm{M}^0) \|=0$, where $\text{HE}(\cdot)$ represents the hyperspherical energy of a weight matrix. For illustration, let us examine a fully connected layer expressed as $\thickmuskip=2mu \medmuskip=2mu \bm{W}=\{\bm{w}_1,\cdots,\bm{w}_n\}\in\mathbb{R}^{d\times n}$, where each neuron $\bm{w}_i\in\mathbb{R}^d$ and $\bm{W}^0$ denotes its pretrained version. The output $\thickmuskip=2mu \medmuskip=2mu \bm{z}\in\mathbb{R}^n$ for input $\thickmuskip=2mu \medmuskip=2mu \bm{x}\in\mathbb{R}^d$ is given by $\thickmuskip=2mu \medmuskip=2mu \bm{z}=\bm{W}^\top\bm{x}$. The OFT method can be interpreted as minimizing the difference in hyperspherical energy between the finetuned and original model weights:

\begin{equation}\label{eq:oft_principle}
\footnotesize
\min_{\bm{W}} \left\| \text{HE}(\bm{W})-\text{HE}(\bm{W}^0)\right\|~~\Leftrightarrow~~\min_{\bm{W}} \bigg{\|} \sum_{i\neq j} \|\hat{\bm{w}}_i-\hat{\bm{w}}_j\|^{-1}-\sum_{i\neq j} \|\hat{\bm{w}}^0_i-\hat{\bm{w}}^0_j\|^{-1}\bigg{\|}
\end{equation}

Here, $\thickmuskip=2mu \medmuskip=2mu \hat{\bm{w}}_i:=\bm{w}_i/\|\bm{w}_i\|$ denotes the normalized $i$-th neuron, and the hyperspherical energy for a weight matrix $\bm{W}$ is defined by $\thickmuskip=2mu \medmuskip=2mu \text{HE}(\bm{W}):= \sum_{i\neq j} \|\hat{\bm{w}}_i-\hat{\bm{w}}_j\|^{-1}$. It is evident that the minimum attainable value for Eq.~\eqref{eq:oft_principle} is zero. This optimum is achieved when $\bm{W}$ can be generated from $\bm{W}^0$ via an orthogonal transformation, meaning $\bm{W}=\bm{R}\bm{W}^0$ for some orthogonal matrix $\bm{R}\in\mathbb{R}^{d\times d}$ (where determinant $1$ and $-1$ indicate rotation and reflection, respectively). The foundational idea behind OFT, then, is that fine-tuning can be accomplished by learning layer-specific orthogonal matrices that rotate or reflect neurons in each layer. In practical terms, OFT optimizes the orthogonal transformation $\thickmuskip=2mu \medmuskip=2mu \bm{R}\in\mathbb{R}^{d \times d

\begin{equation}\label{eq:oft_general}
    \footnotesize
    \bm{z}=\bm{W}^\top\bm{x}=(\bm{R}\cdot\bm{W}^0)^\top\bm{x},~~~\text{s.t.}~\bm{R}^\top\bm{R}=\bm{R}\bm{R}^\top=\bm{I}
\end{equation}

Here, $\bm{W}$ refers to the weight matrix adapted by OFT-finetuning, while $\bm{I}$ is the identity matrix. An illustration of OFT can be found in Figure~\ref{fig:block}. Drawing a parallel to the approach of zero initialization in LoRA, it is crucial for OFT to commence finetuning strictly from the pretrained parameters $\bm{W}^0$. This is accomplished by initializing the orthogonal matrix $\bm{R}$ as an identity matrix, thus ensuring that the starting point of the finetuned model coincides with the pretrained weights. To enforce the orthogonality constraint on $\bm{R}$, one can leverage differentiable orthogonalization techniques such as those presented in \cite{liu2021orthogonal,lezcano2019cheap}. In what follows, we discuss efficient approaches for maintaining this orthogonality property.

\subsection{Efficient Orthogonal Parameterization}\label{sect:eff_ortho}

While traditional orthogonalization methods like the Gram-Schmidt process are differentiable, they tend to be computationally prohibitive for practical applications~\cite{liu2021orthogonal}. To improve efficiency, we utilize the Cayley transform as a means of orthogonal matrix parameterization. In this framework, the orthogonal matrix is realized as $\thickmuskip=2mu \medmuskip=2mu \bm{R}=(\bm{I}+\bm{Q})(I-\bm{Q})^{-1}$, where $\bm{Q}$ is a skew-symmetric matrix satisfying $\thickmuskip=2mu \medmuskip=2mu \bm{Q}=-\bm{Q}^\top$. This efficient method does have a minor trade-off: the Cayley transform exclusively constructs orthogonal matrices with determinant $1$, which form the special orthogonal group. However, our experiments show that this restriction is not detrimental to performance. Despite the efficiency of the Cayley transform, parameterizing a full orthogonal matrix $\bm{R}$ still introduces significant overhead for large $d$. To mitigate this issue, we advocate for representing $\bm{R}$ as a block-diagonal matrix composed of $r$ blocks, leading to the structure:

\begin{equation}
    \footnotesize
    \bm{R}=\text{diag}(\bm{R}_1,\bm{R}_2,\cdots,\bm{R}_r)=\begin{bmatrix}
    \bm{R}_1\in \text{O}(\frac{d}{r}) &  &\\
    &\ddots & \\
    & & \bm{R}_r\in \text{O}(\frac{d}{r})
    \end{bmatrix}\in \text{O}(d)
\end{equation}

where $\text{O}(d)$ stands for the orthogonal group in $d$ dimensions, and $\thickmuskip=2mu \medmuskip=2mu \bm{R}\in\mathbb{R}^{d\times d}$ together with $\thickmuskip=2mu \medmuskip=2mu \bm{R}_i\in\mathbb{R}^{d/r\times d/r},\forall i$ represent orthogonal matrices. If we set $\thickmuskip=2mu \medmuskip=2mu r=1$, then the block-diagonal orthogonal matrix is equivalent to a standard full (unconstrained) orthogonal matrix. For an orthogonal matrix of size $\thickmuskip=2mu \medmuskip=2mu d\times d$, the parameter count is $\thickmuskip=2mu \medmuskip=2mu d(d-1)/2$, giving a parameter complexity of $\mathcal{O}(d^2)$. With $r$-block diagonal structure, the total parameter count reduces to $\thickmuskip=2mu \medmuskip=2mu d(d/r-1)/2$, yielding a complexity of $\thickmuskip=2mu \medmuskip=2mu \mathcal{O}(d^2/r)$. A further reduction is possible by tying the blocks across the matrix, \emph{i.e.}, enforcing $\thickmuskip=2mu \medmuskip=2mu\bm{R}_i=\bm{R}_j$ for all $i\neq j$, which lowers the complexity to $\mathcal{O}(d^2/r^2)$. Notwithstanding these different approaches for compressing parameter count, the final matrix $\bm{R}$ always remains orthogonal, and thus the hyperspherical energy preservation is fully maintained.

We now compare the parameter efficiency between OFT and LoRA. In the case of LoRA, for a specified low-rank parameter $r'$, the number of additional parameters is $\thickmuskip=2mu \medmuskip=2mu r'(d+n)$. Assuming both $r$ and $r'$ scale with $d$, for example, $\thickmuskip=2mu \medmuskip=2mu r=r'=\alpha d$ with some constant $0<\alpha\leq 1$, then LoRA exhibits parameter complexity of $\mathcal{O}(d^2+dn)$ whereas OFT reduces this to $\mathcal{O}(d)$. For illustration, suppose the weight matrix dimension is $\thickmuskip=2mu \medmuskip=2mu 128\times 128$; in this setup, LoRA contains $2{,}048$ trainable parameters when $\thickmuskip=2mu \medmuskip=2mu r'=8$, while OFT contains $960$ trainable parameters with $\thickmuskip=2mu \medmuskip=2mu r=8$ (with block sharing disabled).

\subsection{Constrained Orthogonal Finetuning}

The expressive capacity of standard OFT can be further curtailed by imposing constraints so the finetuned parameters stay close to the pretrained ones. More specifically, COFT employs the following computation in the forward pass:

\begin{equation}\label{eq:coft_general}
    \footnotesize
    \bm{z}=\bm{W}^\top\bm{x}=(\bm{R}\cdot\bm{W}^0)^\top\bm{x},~~~\text{s.t.}~\bm{R}^\top\bm{R}=\bm{R}\bm{R}^\top=\bm{I},~\norm{\bm{R}-\bm{I}}\leq \epsilon
\end{equation}

This formulation places both an orthogonality constraint and an $\epsilon$-proximity condition to the identity. Orthogonality is maintained through the Cayley parameterization introduced in Section~\ref{sect:eff_ortho}. However, imposing the $\epsilon$-bounded deviation in the context of a Cayley-parameterized orthogonal matrix is a challenging problem. To gain further intuition, we use the Neumann series to approximate $\thickmuskip=2mu \medmuskip=2mu \bm{R}=(\bm{I}+\bm{Q})(I-\bm{Q})^{-1}$, obtaining the linear approximation $\thickmuskip=2mu \medmuskip=2mu \bm{R}\

Since the series is convergent in the operator norm, it becomes possible to impose the constraint $\thickmuskip=2mu \medmuskip=2mu\|\bm{R}-\bm{I}\|\leq\epsilon$ directly within the Cayley transform framework. This leads to an equivalent restriction on $\bm{Q}$ given by $\thickmuskip=2mu \medmuskip=2mu\|\bm{Q}-\bm{0}\|\leq\epsilon'$, where $\bm{0}$ stands for the zero matrix and $\epsilon'$ represents a distinct, newly introduced error tolerance parameter. The constraint on $\bm{Q}$ lends itself well to enforcement through projected gradient descent methods. If one wishes to start with $\bm{R}$ set to the identity matrix, initializing $\bm{Q}$ as a zero matrix suffices. The COFT approach, therefore, is characterized by the interplay of two explicit regularizations: maintaining minimal hyperspherical energy variations and enforcing a bounded deviation from the initial model parameters. Although the second type of constraint is often implicitly present in existing finetuning techniques, COFT distinctly formulates it as an explicit regularization term. Notwithstanding OFT's robust empirical performance, we find that the addition of this explicit deviation control in COFT enhances finetuning stability further. Figure~\ref{fig:eps_coft} illustrates the sensitivity of COFT to the parameter $\epsilon$. From these results, it is clear that $\epsilon$ mediates the adaptability allowed during finetuning; increasing $\epsilon$ causes the COFT-finetuned model to increasingly mimic that of OFT, while decreasing $\epsilon$ makes the COFT-finetuned model progressively align with the pretrained diffusion model for text-to-image generation.

\subsection{Re-scaled Orthogonal Finetuning}

We introduce a straightforward augmentation to OFT by permitting a learned scaling factor for each neuron, inspired by the observation that scaling neuron activations does not affect the hyperspherical energy due to subsequent normalization. Concretely, our modified forward computation is given by $\thickmuskip=2mu \medmuskip=2mu\bm{z}=(\bm{R}\bm{W}^0\bm{D})^\top\bm{x}$\footnote{\textbf{Errata}: The NeurIPS camera-ready manuscript mistakenly reports the forward propagation of re-scaled OFT as $\thickmuskip=2mu \medmuskip=2mu\bm{z}=(\bm{D}\bm{R}\bm{W}^0)^\top\bm{x}$. Our implementation, however, follows the correct form, so the empirical findings in Appendix~\ref{app:rescaled} remain valid.}, where $\thickmuskip=2mu \medmuskip=2mu\bm{D}=\text{diag}(s_1,\cdots,s_n)\in\mathbb{R}^{n\times n}$ is a trainable diagonal matrix with strictly positive entries $s_1, \ldots, s_n$. Differing from the initial OFT setup in Eq.~\eqref{eq:oft_general}, where only $\bm{R}$ is optimized, our extension enables both $\bm{D}$ and $\bm{R}$ to be simultaneously updated. This re-scaled OFT method augments OFT's expressive capability while incurring a minimal overhead in parameters. In our reported experiments, we rely on the original OFT framework to focus solely on the effects of orthogonal transformation; however, our investigations indicate that the re-scaled variant consistently achieves stronger performance (see Appendix~\ref{app:rescaled} for details).

\section{Intriguing Insights and Discussions}

\textbf{OFT’s versatility across network architectures}. The OFT technique is, in principle, applicable to any neural architecture. When applied to Transformers, LoRA conventionally operates on attention weights~\cite{hulora2022}, and for fair benchmarking, our experimental protocol restricts OFT-based finetuning exclusively to attention parameters. Beyond standard fully connected layers, OFT readily extends to convolutional networks, exploiting the structure of block-diagonal $\bm{R}$ in meaningful ways for convolution operations—a property not shared by LoRA. If each block in $\bm{R}$ is associated with an individual input channel, every block independently adapts a single neuron channel, reminiscent of depth-wise convolution kernel learning~\cite{chollet2017xception}. In scenarios where a shared block is used across all channels, the same orthogonal transformation is applied to all neurons. We perform an initial exploration of applying OFT for convolution layer finetuning in Appendix~\ref{app:convolution}.

\textbf{Connection to LoRA}. LoRA integrates a low-rank matrix to restrict significant changes to the pretrained weight matrix, thereby preserving existing information. In comparison, OFT enforces the transform applied to the pretrained weight matrix to remain orthogonal and full-rank, which similarly safeguards the pretraining knowledge from being degraded. The forward computation in OFT can be expressed as $\thickmuskip=2mu \medmuskip=2mu \bm{z}=(\bm{R}\bm{W}^0)^\top\bm{x}=(\bm{W}^0+(\bm{R}-\bm{I})\bm{W}^0)^\top\bm{x}$, highlighting that $\thickmuskip=2mu \medmuskip=2mu (\bm{R}-\bm{I})\bm{W}^0$ plays a role comparable to LoRA’s low-rank weight modification. Since the pretrained matrix $\bm{W}^0$ is usually full-rank, OFT achieves a low-rank update when $\thickmuskip=2mu \medmuskip=2mu \bm{R}-\bm{I}$ itself is low-rank. Much like LoRA utilizes a rank parameter $r'$ to limit the parameter count, OFT features a diagonal block parameter $r$ to reduce the number of adjustable parameters. Notably, LoRA and OFT exemplify two contrasting approaches to parameter efficiency: LoRA leverages a low-rank configuration to minimize trainable variables, whereas OFT adopts a block-diagonal, sparse orthogonal structure to achieve similar efficiency gains.

\textbf{Why OFT converges faster}. Evidence in Figure~\ref{fig:toy_ae} indicates that the most impactful way to alter a model's semantics is by modifying the directions of its neurons, a capability for which OFT is specifically developed. Furthermore, OFT can be interpreted as adapting the neurons along a smooth hypersphere manifold, a process which improves the optimization properties. This benefit has also been demonstrated empirically in \cite{liu2021orthogonal}.

\textbf{Why not minimize hyperspherical energy}. Unlike \cite{liu2021orthogonal}, our method does not attempt to reduce hyperspherical energy. For classification tasks, a non-redundant neuronal arrangement is typically preferred. Achieving minimal hyperspherical energy would result in neurons being evenly distributed on the hypersphere, which is not a suitable goal during finetuning since it risks erasing the useful information encoded during pretraining.

\textbf{Balancing flexibility and regularity in finetuning}. Our investigation reveals an intrinsic compromise between flexibility and regularity. While conventional finetuning techniques exhibit maximal adaptability, they often suffer from instability and are prone to model degeneration. OFT, despite its conceptual simplicity, effectively strikes a balance by maintaining pairwise neuron angular relationships, ensuring both adaptability and structural discipline. The block-diagonal parameterization further acts as a robust form of regularization on the orthogonal matrix.

\textbf{No inference-time computational cost}. In contrast to approaches like ControlNet, the OFT method does not introduce any additional computational burden during inference. At test time, the finetuned model is constructed by multiplying the optimized orthogonal matrix $\bm{R}$ with the original pretrained weight matrix $\bm{W}^0$, resulting in $\thickmuskip=2mu \medmuskip=2mu \bm{W}=\bm{R}\bm{W}^0$. Consequently, inference throughput remains unchanged relative to the base model.

\section{Experiments and Results}

\textbf{Experimental setup}. For all experiments, Stable Diffusion v1.5~\cite{rombach2022high} serves as the foundational text-to-image backbone. Image outputs from each approach are selected at random to preserve impartiality. The procedure for subject-driven synthesis largely follows DreamBooth~\cite{ruiz2023dreambooth}, while protocols for controllable image generation adhere to guidelines from ControlNet~\cite{zhang2023adding} and T2I-Adapter~\cite{mou2023t2i}. For a consistent comparison with LoRA, OFT or COFT modifications are limited to the same architectural layers as those modified by LoRA. Expanded experimental results and further procedural clarification appear in Appendix~\ref{app:settings}.

\subsection{Subject-driven Generation}

\textbf{Configurations}. DreamBooth~\cite{ruiz2023dreambooth} and LoRA~\cite{hulora2022} are selected as baseline models. All evaluated techniques leverage the loss formulation introduced by DreamBooth. Hyperparameter choices for DreamBooth and LoRA are aligned with their respective original publications, employing their optimal configurations. Comprehensive outcomes and supplementary experimental insights can be found in Appendix~\ref{app:settings},\ref{app:coft_oft},\ref{app:more_qualitative},\ref{app:failure}.

\textbf{Finetuning stability and convergence}. We begin by assessing the stability during finetuning and the rate of convergence for DreamBooth, LoRA, OFT, and COFT. Figure~\ref{fig:teaser} and Figure~\ref{fig:db_convergence} present these results. It is evident that both COFT and OFT deliver reliable and steady finetuning of the diffusion model. After 400 steps, DreamBooth and OFT-based methods exert effective control, yet LoRA fails to maintain the subject’s identity. After progressing to 2000 iterations, DreamBooth exhibits mode collapse by producing degenerate images, and LoRA is unable to correctly represent the yellow shirt, instead yielding yellow fur. Conversely, both OFT and COFT persist in providing stable and robust control over image generation. These findings demonstrate the efficiency and stability of the OFT approach for subject-driven generation tasks. Notably, the robustness of OFT and COFT to the chosen number of finetuning steps is especially valuable, as it relieves practitioners from the burden of meticulously selecting iteration counts for different subjects. With OFT and COFT, one may simply choose a sufficiently large number of iterations without tuning. Furthermore, in the case of COFT, a suitable selection of $\epsilon$ enables both the learning rate and number of steps to be set with minimal effort.

\textbf{Quantitative comparison}. Adopting the evaluation protocols from \cite{ruiz2023dreambooth}, we perform quantitative assessments to measure subject fidelity (DINO~\cite{caron2021emerging}, CLIP-I~\cite{radford2021learning}), alignment to text prompts (CLIP-T~\cite{radford2021learning}), and diversity among generated samples (LPIPS~\cite{zhang2018unreasonable}). In particular, CLIP-I is computed by averaging the pairwise cosine similarity of CLIP embeddings between the real and synthesized images, while DINO follows the same computation using ViT S/16 DINO features. CLIP-T captures the mean cosine similarity between the text prompts and their respective generated images, as measured by CLIP embeddings. For diversity, the average LPIPS cosine similarity is calculated across images produced from identical prompts for the same subject. The scores in Table~\ref{table:dreambooth_final} indicate that COFT and OFT both markedly surpass DreamBooth and LoRA on the DINO and CLIP-I subject fidelity measures, while also achieving slightly superior or on-par results in terms of prompt consistency and image diversity. To control for randomness, each model is finetuned 30 times from the same pretrained checkpoint, each time with a different random seed. Collectively, these results substantiate that OFT not only enhances convergence stability but also secures consistently stronger final outcomes.

\textbf{Qualitative comparison}. To provide clearer insight into the advantages of OFT, we present several randomly selected cases of subject-driven generation in Figure~\ref{fig:dreambooth_final}. For consistency and fairness, every example shown is produced using the same finetuned model across all methods, with no prompt-specific optimization carried out for any final output. We evaluate each method based on the model that attains the highest validation CLIP score. Examining Figure~\ref{fig:dreambooth_final}, it becomes apparent that both OFT and COFT achieve robust semantic subject retention, whereas LoRA frequently fails to maintain subject identity (\emph{e.g.}, in the bowl example, LoRA does not preserve the subject’s features at all). Additionally, both OFT and COFT exhibit superior text-driven control in image generation; in contrast, DreamBooth, though effective in retaining subject identity, frequently does not adhere to the descriptive prompt (\emph{e.g.}, see the first row for the bowl scenario). These qualitative results collectively indicate that the OFT framework excels at balancing image controllability and faithful subject preservation. Furthermore, OFT demonstrates stability with respect to the number of training iterations, maintaining strong performance even as the iteration count increases—a robustness not observed in either DreamBooth or LoRA. Supplementary qualitative examples can be found in Appendix~\ref{app:more_qualitative}. In addition, results from a human study, detailed in Appendix~\ref{app:human}, offer additional confirmation of OFT's superior performance.

\subsection{Controllable Generation}

\textbf{Settings}. As comparative baselines, we include ControlNet~\cite{zhang2023adding}, T2I-Adapter~\cite{mou2023t2i}, and LoRA~\cite{hulora2022}. The main paper addresses three demanding controllable generation benchmarks: translating Canny edge maps to images~(C2I) using the COCO dataset~\cite{lin2014microsoft}, converting segmentation maps to images~(S2I) based on the ADE20K dataset~\cite{zhou2017scene}, and generating faces from landmarks~(L2F) with CelebA-HQ~\cite{karras2018progressive,xia2021tedigan}. All approaches involve fine-tuning the Stable Diffusion~(SD) v1.5 model on these datasets for a total of 20 training epochs. Further outcomes can be found in Appendix~\ref{app:more_qualitative},~\ref{app:more_control_task},~\ref{app:failure}.

\textbf{Convergence}. To assess convergence speed, we compare ControlNet, T2I-Adapter, LoRA, and COFT within the L2F task, using both quantitative and qualitative measures. For the quantitative analysis, our metric is the average $\ell_2$ distance calculated between the set of control face landmarks and their predicted counterparts. Figure~\ref{fig:face_convergence} depicts the face landmark errors achieved by models fine-tuned across varying numbers of epochs. From these results, it is clear that both COFT and OFT attain much faster convergence than alternatives. For instance, LoRA requires 20 epochs to reach the level of performance that our OFT achieves by the 8th epoch. Importantly, the trainable parameter counts of OFT and COFT are on par with that of LoRA (and are substantially less than those of ControlNet), yet they demonstrate superior efficiency in terms of convergence. The accelerated learning of OFT is further confirmed by Figure~\ref{fig:teaser}, where the example on the right illustrates that OFT exhibits much greater data efficiency than either ControlNet or LoRA, being capable of successful convergence with just 5\% of the ADE20K dataset. For qualitative comparisons, we focus on OFT, COFT, and ControlNet, given that ControlNet attains the most comparable landmark error to ours. As demonstrated in Figure~\ref{fig:face_convergence_qual}, both OFT and COFT display consistent and stable convergence, with the synthesized face pose incrementally aligning with the control landmarks. By contrast, ControlNet shows a sudden gain in controllability after the 8th epoch, echoing the abrupt convergence behavior documented in \cite{zhang2023adding}. The steady and predictable convergence of OFT greatly facilitates practical usage, as its training trajectory is considerably smoother and easier to anticipate, which simplifies the search for effective OFT hyperparameters.

\textbf{Quantitative comparison}. To assess the efficacy of controllable generation, we propose a metric that quantifies control consistency. This approach involves extracting the control signal from the synthesized image and measuring its similarity to the original control input. Specifically, for the C2I task, we report both IoU and F1 scores. For S2I, we evaluate using mean IoU as well as mean and overall accuracy metrics. In the L2F task, we use the mean $\ell_2$ distance between the control landmarks and those predicted. Comprehensive descriptions of these consistency metrics can be found in Appendix~\ref{app:settings}. All methods are evaluated under their respective optimal hyperparameter configurations. According to Table~\ref{table:control_gen}, OFT and COFT consistently demonstrate superior and more precise controllability relative to other evaluated techniques. It is apparent that adapter-based solutions (\emph{e.g.}, T2I-Adapter and ControlNet) not only exhibit slower convergence but also tend to deliver inferior results. In a comparison of performance, LoRA surpasses ControlNet on the S2I task yet falls short in both C2I and L2F scenarios. Notably, we observe that LoRA’s ultimate effectiveness appears limited, even under extensive hyperparameter optimization. By contrast, the OFT framework does not exhibit performance saturation and empirically continues to improve as its trainable parameter set expands. We highlight that our introduced quantitative protocol for assessing controllable generation constitutes a key innovation, as it enables robust evaluation of finetuned model control fidelity, constrained only by the baseline accuracy of existing segmentation or detection models.

\textbf{Qualitative comparison}. We further conduct a qualitative assessment of OFT and COFT in relation to leading approaches such as ControlNet, T2I-Adapter, and LoRA. As illustrated in Figure~\ref{fig:control_final}, randomly sampled results indicate that OFT and COFT generate images of high realism and fidelity, while also enabling precise control over the outputs. Specifically, for the S2I task, LoRA is unable to synthesize images consistent with the provided segmentation maps, whereas ControlNet, OFT, and COFT effectively follow the given input for image generation. Unlike ControlNet, OFT and COFT are capable of producing high-quality, detail-rich, and geometrically plausible images, all with substantially fewer model parameters. In the C2I setting, both OFT and COFT can convincingly create semantically aligned images from sparse Canny edge maps, outperforming T2I-Adapter and LoRA, which show considerably poorer results. When it comes to the L2F task, our approach exhibits the highest accuracy in maintaining pose control for generated faces, even in cases of complex head orientations. Across the three examined control tasks, OFT and COFT consistently yield higher qualitative performance compared to current baselines, attesting to the robustness of our OFT framework for controllable image generation. To support these findings, we present additional qualitative samples for all three tasks in Appendix~\ref{app:control_exp}, and further show the versatility of OFT across extra control scenarios (including dense pose to human body, sketch to image, and depth to image tasks) in Appendix~\ref{app:more_control_task}.

\section{Concluding Remarks and Open Problems}

Inspired by the insight that the angular relationships between neurons play a vital role in defining visual semantics, we introduce a straightforward yet powerful finetuning technique—orthogonal finetuning—for regulating text-to-image diffusion models. This approach is particularly tailored for two scenarios: subject-driven generation and controllable generation. OFT, when compared to previous methods, achieves enhanced controllability and greater stability during finetuning, all while utilizing a reduced set of finetuning parameters. Notably, OFT does not cause any extra inference latency, making it suitable for practical deployment.

However, the adoption of OFT raises several noteworthy open questions. Firstly, OFT ensures orthogonality through Cayley parametrization, which necessitates computing a matrix inverse and thereby constrains the scalability of the approach. While we mitigate this bottleneck using block diagonal parametrization, devising faster, differentiable matrix inverse computations remains an unresolved issue. Secondly, the structure of OFT presents promising opportunities for compositionality: the orthogonal matrices yielded from various OFT finetuning processes can be multiplied together, preserving their orthogonality. An intriguing avenue for future investigation is whether such a collection of orthogonal matrices can effectively retain the task-specific information from all downstream tasks. Lastly, OFT’s parameter efficiency is primarily determined by its block diagonal design, a factor that may inherently introduce bias and limit adaptability. Developing more efficient and less biased strategies for parameter utilization is thus a key area for further research.

\end{document}